\documentclass[10pt]{beamer}
\usetheme{metropolis}
\usepackage{graphicx}
\usepackage{listings}
\usepackage{booktabs}
\usepackage{xcolor}
\definecolor{codebg}{HTML}{F5F5F5}
\definecolor{codegreen}{HTML}{2E7D32}
\definecolor{codegray}{HTML}{757575}
\definecolor{codeblue}{HTML}{1565C0}
\lstdefinestyle{rstyle}{language=R,backgroundcolor=\color{codebg},basicstyle=\ttfamily\scriptsize,
  keywordstyle=\color{codeblue}\bfseries,stringstyle=\color{codegreen},
  commentstyle=\color{codegray}\itshape,breaklines=true,frame=single,
  rulecolor=\color{codegray!30},showstringspaces=false,tabsize=2,
  morekeywords={library,ggplot,aes,geom_point,geom_line,geom_col,geom_histogram,geom_boxplot,
    facet_wrap,facet_grid,patchwork,plot_annotation,tag_levels,plot_layout,
    theme_minimal,theme,element_text,labs,wrap_plots,free}}
\lstdefinestyle{pythonstyle}{language=Python,backgroundcolor=\color{codebg},basicstyle=\ttfamily\scriptsize,
  keywordstyle=\color{codeblue}\bfseries,stringstyle=\color{codegreen},
  commentstyle=\color{codegray}\itshape,breaklines=true,frame=single,
  rulecolor=\color{codegray!30},showstringspaces=false,tabsize=4,
  morekeywords={import,from,as,pd,plt,sns,np,fig,ax,subplots,GridSpec,
    add_subplot,savefig,suptitle,tight_layout}}
\title{Small Multiples \& Faceted Displays}
\subtitle{Workshop 3 --- Module 8}
\author{Prof.Asoc.\ Endri Raco}
\institute{Data Visualization for Data Scientists \\ UNYT Tirana}
\date{2025/26}
\begin{document}

\begin{frame}
  \titlepage
\end{frame}

\begin{frame}{Module Roadmap}
  \begin{enumerate}
    \item Tufte's small multiples revisited as a chart type
    \item Cleveland's Trellis philosophy: four principles
    \item Composition in R: patchwork operators and cowplot
    \item Composition in Python: GridSpec layout patterns
    \item Facet (homogeneous) vs patchwork (heterogeneous) decision
    \item Building multi-panel dashboards as static figures
  \end{enumerate}
  \begin{exampleblock}{Learning Outcome}
    You will compose heterogeneous multi-panel figures using
    \texttt{patchwork} (R) and \texttt{GridSpec} (Python), choosing
    between faceting and manual composition based on panel homogeneity.
  \end{exampleblock}
\end{frame}

\section{From Spaghetti to Small Multiples}

\begin{frame}{The Motivation: 6 Lines $\rightarrow$ 6 Panels}
  \begin{center}
    \includegraphics[width=0.95\textwidth]{figures_w03m08/spaghetti_to_sm}
  \end{center}
\end{frame}

\begin{frame}{Trellis Philosophy: Four Principles}
  \begin{center}
    \includegraphics[width=0.92\textwidth]{figures_w03m08/trellis_philosophy}
  \end{center}
  \begin{alertblock}{Pro-Tip}
    Cleveland (1993) formalised small multiples as \textbf{Trellis
    displays}: same chart type, same axes, one variable per panel.
    This exploits the pre-attentive system: the reader scans panels
    by position (Gestalt proximity from W01-M03) rather than by colour
    discrimination.
  \end{alertblock}
\end{frame}

\section{Composition in R: patchwork}

\begin{frame}{patchwork Operators}
  \begin{center}
    \includegraphics[width=0.85\textwidth]{figures_w03m08/patchwork_operators}
  \end{center}
\end{frame}

\begin{frame}[fragile]{patchwork: Complete Example}
\begin{lstlisting}[style=rstyle]
library(patchwork)

p_bar <- ggplot(df, aes(x = fct_reorder(cat, n), y = n)) +
  geom_col(fill = "#1565C0") + coord_flip() + theme_minimal()

p_scatter <- ggplot(df, aes(x = reviews, y = rating)) +
  geom_point(alpha = 0.3) + scale_x_log10() + theme_minimal()

p_hist <- ggplot(df, aes(x = rating)) +
  geom_histogram(bins = 30, fill = "#2E7D32") + theme_minimal()

p_box <- ggplot(df, aes(x = type, y = rating)) +
  geom_boxplot(fill = "#E3F2FD") + theme_minimal()

# Compose: | = side by side, / = stack
dashboard <- (p_bar | p_scatter) / (p_hist | p_box) +
  plot_annotation(
    title = "Google Play Store: Visual EDA",
    tag_levels = "a") &      # & applies to ALL panels
  theme_minimal(base_size = 9)

ggsave("dashboard.pdf", dashboard, width = 12, height = 8)
\end{lstlisting}
\end{frame}

\begin{frame}[fragile]{patchwork: Advanced Layouts}
\begin{lstlisting}[style=rstyle]
# Unequal widths: bar gets 2x space
(p_bar + p_scatter) + plot_layout(widths = c(2, 1))

# Nested: bar on left, scatter + hist stacked on right
p_bar | (p_scatter / p_hist)

# Insets: small plot inside a large one
p_scatter + inset_element(p_hist, left = 0.6, bottom = 0.6,
                          right = 1, top = 1)

# cowplot alternative (for legacy code)
library(cowplot)
plot_grid(p_bar, p_scatter, p_hist, p_box,
          ncol = 2, labels = "auto",
          rel_widths = c(1.5, 1))
\end{lstlisting}
  \begin{alertblock}{Pro-Tip}
    \texttt{patchwork} is the modern standard (Thomas Lin Pedersen, 2020).
    \texttt{cowplot} (Claus Wilke) and \texttt{gridExtra} are older
    alternatives still found in legacy code. Prefer patchwork for new work.
  \end{alertblock}
\end{frame}

\section{Composition in Python: GridSpec}

\begin{frame}{GridSpec: Three Layout Patterns}
  \begin{center}
    \includegraphics[width=0.92\textwidth]{figures_w03m08/gridspec_layouts}
  \end{center}
\end{frame}

\begin{frame}[fragile]{GridSpec: Complete Example}
\begin{lstlisting}[style=pythonstyle]
from matplotlib.gridspec import GridSpec

fig = plt.figure(figsize=(14, 9))
gs = GridSpec(2, 3, figure=fig, hspace=0.4, wspace=0.35)

# Top row: 3 equal panels
ax_bar  = fig.add_subplot(gs[0, 0])    # (a)
ax_hist = fig.add_subplot(gs[0, 1])    # (b)
ax_scat = fig.add_subplot(gs[0, 2])    # (c)

# Bottom row: 1 wide + 2 narrow
ax_line = fig.add_subplot(gs[1, :2])   # (d) spans 2 columns
ax_box  = fig.add_subplot(gs[1, 2])    # (e) single column

# ... plot each panel onto its axes ...

# Global title + caption
fig.suptitle("Dashboard Title", fontsize=13, fontweight="bold")
fig.text(0.5, 0.01, "Source: ...", ha="center", fontsize=8, color="#888")

fig.savefig("dashboard.pdf", bbox_inches="tight")
\end{lstlisting}
\end{frame}

\begin{frame}[fragile]{subplot\_mosaic: Named Layouts (Python 3.4+)}
\begin{lstlisting}[style=pythonstyle]
# matplotlib 3.4+ : named subplot layout
fig, axes = plt.subplot_mosaic(
    [["bar", "bar", "scatter"],
     ["hist", "box", "scatter"]],
    figsize=(12, 8))

# Access by name instead of index
axes["bar"].barh(cats, vals, color="#1565C0")
axes["scatter"].scatter(x, y, s=10, alpha=0.3)
axes["hist"].hist(ratings, bins=30, color="#2E7D32")
axes["box"].boxplot([free_r, paid_r], labels=["Free","Paid"])
\end{lstlisting}
  \begin{exampleblock}{Data Insight}
    \texttt{subplot\_mosaic} uses an ASCII grid to define layout:
    repeated names span multiple cells (like \texttt{gs[:2]} but more
    readable). Available since matplotlib 3.4. For older versions,
    use \texttt{GridSpec}.
  \end{exampleblock}
\end{frame}

\section{Facet vs Patchwork}

\begin{frame}{When to Facet vs When to Compose}
  \begin{center}
    \includegraphics[width=0.92\textwidth]{figures_w03m08/facet_vs_patchwork}
  \end{center}
\end{frame}

\begin{frame}{Decision Rule}
  \centering
  \begin{tabular}{lll}
    \toprule
    \textbf{Question} & \textbf{Answer} & \textbf{Tool} \\
    \midrule
    Same geom in all panels? & Yes & \texttt{facet\_wrap/grid} \\
    Same aes in all panels? & Yes & \texttt{facet\_wrap/grid} \\
    Different geoms per panel? & Yes & \texttt{patchwork / GridSpec} \\
    Need unequal panel sizes? & Yes & \texttt{patchwork(widths=)} \\
    Need inset plot? & Yes & \texttt{patchwork::inset\_element()} \\
    Need panel labels (a,b,c)? & Yes & \texttt{tag\_levels = "a"} \\
    \bottomrule
  \end{tabular}

  \vspace{6pt}
  {\small The simplest rule: if you can write one \texttt{ggplot() + facet\_*()}
  call, use faceting. If you need multiple separate \texttt{ggplot()} calls
  with different geoms, use patchwork.}
\end{frame}

\section{Dashboard Example}

\begin{frame}{Six-Panel Dashboard via GridSpec}
  \begin{center}
    \includegraphics[width=0.88\textwidth]{figures_w03m08/dashboard}
  \end{center}
\end{frame}

\section{Complete Examples}

\begin{frame}[fragile]{R: patchwork 2$\times$2}
  \begin{columns}[T]
    \begin{column}{0.52\textwidth}
\begin{lstlisting}[style=rstyle]
library(patchwork)

p1 <- ggplot(df, aes(x=cat,y=n)) +
  geom_col(fill="#1565C0") +
  coord_flip() + theme_minimal() +
  labs(title="(a) Bar")

p2 <- ggplot(df, aes(x=x,y=y)) +
  geom_point(alpha=0.5, size=1,
    color="#E53935") +
  theme_minimal() +
  labs(title="(b) Scatter")

p3 <- ggplot(df, aes(x=val)) +
  geom_histogram(bins=20,
    fill="#2E7D32") +
  theme_minimal() +
  labs(title="(c) Histogram")

p4 <- ggplot(df, aes(x=1:10,y=y)) +
  geom_line(color="#7B1FA2") +
  theme_minimal() +
  labs(title="(d) Line")

(p1 | p2) / (p3 | p4)
\end{lstlisting}
    \end{column}
    \begin{column}{0.45\textwidth}
      \includegraphics[width=\textwidth]{figures_w03m08/r_output}
    \end{column}
  \end{columns}
\end{frame}

\begin{frame}[fragile]{Python: GridSpec Asymmetric Layout}
  \begin{columns}[T]
    \begin{column}{0.52\textwidth}
\begin{lstlisting}[style=pythonstyle]
fig = plt.figure(figsize=(10, 6))
gs = GridSpec(2, 3, figure=fig,
    hspace=0.4, wspace=0.3)

# (a) Top wide bar (spans all 3 cols)
ax_bar = fig.add_subplot(gs[0, :])
ax_bar.barh(cats, vals,
    color="#1565C0", height=0.5)
ax_bar.set_title("(a) Wide bar")

# (b,c,d) Bottom row: 3 panels
ax_sc = fig.add_subplot(gs[1, 0])
ax_sc.scatter(x, y, s=8,
    c="#E53935", alpha=0.5)
ax_sc.set_title("(b) Scatter")

ax_hi = fig.add_subplot(gs[1, 1])
ax_hi.hist(data, bins=15,
    color="#2E7D32")
ax_hi.set_title("(c) Hist")

ax_bx = fig.add_subplot(gs[1, 2])
ax_bx.boxplot([d1], widths=0.4,
    patch_artist=True)
ax_bx.set_title("(d) Box")
\end{lstlisting}
    \end{column}
    \begin{column}{0.45\textwidth}
      \includegraphics[width=\textwidth]{figures_w03m08/py_output}
    \end{column}
  \end{columns}
\end{frame}

\section{Synthesis}

\begin{frame}{Key Takeaways}
  \begin{enumerate}
    \item \textbf{Small multiples} (Tufte) = same chart, same axes,
          one variable per panel. Exploits position over colour.
    \item \textbf{Trellis principles} (Cleveland): same geom, shared axes,
          one variable per panel, minimal strip labels.
    \item \textbf{patchwork} (R): \texttt{|} = side by side, \texttt{/} =
          stack, \texttt{\&} = apply to all, \texttt{tag\_levels} = auto
          panel labels.
    \item \textbf{GridSpec} (Python): \texttt{gs[row, col]} for placement,
          \texttt{gs[0, :]} for spanning. \texttt{subplot\_mosaic} for
          named ASCII layouts.
    \item Use \textbf{faceting} for homogeneous panels (same geom/aes);
          \textbf{patchwork/GridSpec} for heterogeneous dashboards.
    \item Every multi-panel figure needs: \textbf{panel labels} (a, b, c),
          \textbf{shared title}, and \textbf{source caption}.
  \end{enumerate}
\end{frame}

\begin{frame}{Essential Readings}
  \begin{itemize}
    \item Cleveland, W.~S. (1993). \emph{Visualizing Data}. (Trellis displays.)
    \item Tufte, E.~R. (1983). \emph{The Visual Display of Quantitative
          Information}, pp.\ 170--174 (Small multiples.)
    \item Pedersen, T.~L. (2020). patchwork: The Composer of Plots.
          \texttt{https://patchwork.data-imaginist.com}
    \item Wilke, C.~O. (2019). \emph{Fundamentals of Data Visualization},
          Chapter 21 (Multi-Panel Figures).
  \end{itemize}
  \vspace{6pt}
  \textbf{Next Module}: Tables as Visualization (W03-M09)
\end{frame}

\end{document}
